\documentclass[11pt]{article}
\usepackage{fontspec}
\setmainfont{Times New Roman}
\setsansfont{Arial}
\setmonofont{Consolas}
\usepackage{amsmath,amssymb,mathtools}
\usepackage{geometry}
\usepackage{microtype}
\newcommand{\mS}{\mathfrak{S}}
\newcommand{\mA}{\mathfrak{A}}
\newcommand{\mo}{\mathfrak{o}}
\newcommand{\mO}{\mathfrak{O}}
\newcommand{\mT}{\mathfrak{T}}
\newcommand{\mR}{\mathfrak{R}}
\newcommand{\mq}{\mathfrak{q}}
\geometry{a4paper,margin=2.35cm}
\setlength{\parindent}{1.25em}
\setlength{\parskip}{0pt}
\makeatletter
\renewcommand\footnotesize{\@setfontsize\footnotesize{7.6}{8.6}\selectfont}
\makeatother
\emergencystretch=100em
\allowdisplaybreaks
\sloppy
\hbadness=10000
\vbadness=10000
\hfuzz=2pt
\vfuzz=10pt
\raggedbottom

\begin{document}

% Unit: Paper32 section 1 reduction theorem v001. Direct Ukrainian translation from the R122 German Paper32/Paper33 source slice, lines 380 and 382 / segments P32-S0008--P32-S0009.
\noindent Нехай тепер $\mS$ припускається цілком звідною над $P_0$; нехай $P_0$ вважається досконалим полем, так що те саме виконується для кожного алгебричного розширення $P$ поля $P_0$. Класи зображень за кожного розширення $P$ поля $P_0$ залишаються цілком звідними.\footnote{Якщо $P_0$ не є досконалим, цілковита звідність зберігається тоді й лише тоді, коли компоненти $K$ центра стають “розширеннями першого роду” поля $P_0$. За цієї умови чинними лишаються всі дальші наслідки тексту.} Центр $\mS$ є прямою сумою скінченної кількості полів; кожне розширення $P$ поля $P_0$, ізоморфне одному з таких полів $K$, є полем одного (простого) характеру. Коли розширити $P_0$ до ізоморфного $K$ поля $P$, від $\mS$ відщеплюється двобічно просте кільце, якому відповідає абсолютно незвідний клас зображень цього характеру. Це кільце -- за теоремою Веддерберна -- є прямим добутком матричного кільця над $P$ із некомутативним тілом $\mA$, однозначно визначеним з точністю до ізоморфізму; центр $\mA$ дорівнює $P$, а його ступінь над $P$ дорівнює $m^2$, де $m$ позначає індекс Шура, що належить цьому характеру. Кожне поле розщеплення $\mS$ -- щодо зображення, яке належить цьому характеру -- водночас є таким для $\mA$ щодо зображень над $P$, і навпаки. Спряжені зображення $\mS$ отримують, виходячи з відповідного двобічно простого кільця над $P_0$; його відповідне тіло $\mA_0$ стає ізоморфним до $\mA$ і має центр $K$. Отже, задачу про поля розщеплення повністю зведено до задачі про відповідне некомутативне тіло $\mA$ та його розщеплення.\footnote{Це відповідає Шуровому зведенню до системи матриць над $P$, які перестановні з незвідним зображенням $\mS$ над $P$. Транспоновані цих матриць дають зображення $\mA$ над $P$.} Зокрема, маємо теорему:

\medskip

\noindent Усі найбільші комутативні підполя $\mA$ мають однаковий ступінь $m$ над $P$, де $m$ означає індекс Шура. Усі ці найбільші комутативні підполя є полями розщеплення найменшого ступеня, і кожне поле розщеплення найменшого ступеня міститься серед них. Ступінь кожного поля розщеплення є кратним $m$. $\mA$ має лише один абсолютно незвідний клас зображень, а також лише один незвідний клас зображень над $P$.

\end{document}
